\section{作業段の仮定を残す}
\label{sec:contract}

地位は\dstopen{}である。Lean に積分器は無い。この章の一段は、実行してよい作業手順であり、論文の運動方程式ではない。仮定をログに残さない呼び出しは、ライブラリが拒否する。

\subsection{一段が読んでよいもの}

\begin{enumerate}
\item 位置はチャート \(E\) の \texttt{ChartBody.pos} だけである。
\item ペアの力は第\ref{sec:force}章の静的な和である。重力和と電磁和を両方使うなら、足し算をこの作業段で行い、その旨をログに書く。
\item 内在角は第\ref{sec:written}章の閉形式である。剛性 \(m\) は \texttt{WrittenSeed.stiffness} として外から渡す。層力の大きさから \(m\) を更新しない。
\item 位置を動かすときに加速度 \(\ddot x=F/\mu\) を使うなら、\(\mu\) は粒子ごとの慣性パラメータである。剛性 \(m\) でも符号なし質量 \(M\) でもない。
\end{enumerate}

\subsection{ログ}

各段の前に、次のレコードを残す。

\begin{verbatim}
struct StepContract:
    tag: "working-step"
    position_law: "frozen_accel_taylor"
    mu: map iota -> real
    stiffness_m: real
    killing_M: real
    used_gravity_sum: bool
    used_em_sum: bool
    note: "each mu_i, stiffness_m, and killing_M stay separate fields"
\end{verbatim}

\(\mu_i\)、\(m\)、\(M\) の数値が偶然一致しても、フィールドは分けたまま残す。一致を恒等式として簡約しない。

\subsection{一段の手順}

\begin{verbatim}
function working_step(bodies, seeds, mu, dt, inputs) -> reject or state:
    require step contract log
    for each body i:
        F_i = gravityLayerForce on the other bodies
        if em requested: F_i = F_i + em sum on the other bodies
        if any pair has r = 0 or gamma = 0: reject before writing
    for each body i:
        a = F_i / mu_i
        pos_i = pos_i + vel_i * dt + a * dt * dt / 2
        vel_i = vel_i + a * dt
    for each seed:
        sample written_at(seed, t + dt)
        build alpha, beta
        if not admissible: reject
    t = t + dt
\end{verbatim}

位置の差分は、段のあいだ加速度を一定とみなしたテイラー更新の例である。DST の定理ではない。別の段法を使うときは \texttt{position\_law} の文字列を変える。内在角の側は、段法を替えても閉形式のままにする。閉形式は錐の外でも式として定義されている。錐の検査は、シミュレーターが状態として保持するかどうかを決める門である。

\subsection{論文へ昇格させない}

この一段を、日本語論文の「書かれた作用」の節へ書き戻さない。Lean に \texttt{working\_step} という定義を足さない。多体系を閉じる式が足りない、という事実は、この仮定のログが引き受ける。
